\section{INTRODUCTION}
\label{sec:introduction}


The Cosmic Microwave Background (CMB) provides a unique window to probe the early stages of the Universe and its structure formation. Furthermore, CMB polarization enables to probe the amplitude of primordial Gravitational Waves (GWs; \citealp{BICEP2}) generated during the inflationary epoch \citep{CMB_tests_of_inflation, CMB_polarization}. These GWs give rise to the so-called B-mode of CMB polarization. In recent years, many experiments have been observing the CMB anisotropies. For instance, the Wilkinson Microwave Anisotropy Probe (WMAP; \citealp{WMAP_9_year_results}), scanning the full sky in five frequency bands between 23 and 94 GHz, and the Planck satellite \citep{Planck_2018VI} in 9 bands, with polarization measurements in 7 of them, spanning from 30 to 353 GHz. These observations have significantly advanced our understanding of the Universe and motivated further scientific exploration, marking the beginning of a new era of high-precision cosmology with future missions centred in the study of the CMB, such as LiteBIRD \citep{LiteBIRD} or other cosmology experiments such as the Square Kilometre Array (SKA; \citealp{ska}), offering enhanced sensitivity and resolution.

\paragraph{}

The primordial B-mode signal might have an amplitude high enough to be within reach of current and upcoming experiments. However, its detection does not depend solely on instrumental sensitivity. A successful discovery critically relies on accurate component separation, since the cosmological signal is contaminated by polarized foreground emission from the Galaxy with significantly higher amplitude \citep{BICEP2, Planck_2015XIX}. Component separation is therefore a central challenge in the design of future CMB experiments such as LiteBIRD \citep{LiteBIRD}, directly linking the search for primordial B modes to the statistical characterization and astrophysical understanding of polarized emission from the magnetized interstellar medium (ISM).

\paragraph{}

Microwave observations also contain several astrophysical emission mechanisms other than the CMB, collectively known as foregrounds. Because the primordial B-mode signal has a very low amplitude, the removal of these foregrounds in polarization has become a crucial step for its detection. Two foreground components are known to produce strongly linearly polarized radiation: synchrotron emission, generated by relativistic cosmic-ray electrons spiraling around Galactic magnetic field (GMF) lines and dominating at low frequencies; and thermal dust emission from interstellar dust grains aligned with the magnetic field, which dominates at high frequencies. Both components can exhibit polarization fractions of up to $\sim$20$\%$ in certain regions of the sky \citep{Planck_2015XIX}.

Anomalous microwave emission (AME), on the other hand, is observed to have a very low polarization fraction \citep{QUIJOTE_II}. If AME originates from spinning dust grains, theoretical models predict negligible polarization levels \citep{AME_polarization, AME_pol_limits}. Another relevant foreground is free-free emission, produced by free electrons interacting with ionized gas; however, this mechanism generates radiation that is effectively unpolarized \citep{WMAP_7yr_foregrounds}.


\paragraph{}

Synchrotron emission accounts for most of the radio emission from Active Galactic Nuclei (AGNs), which are thought to be powered by supermassive black holes in galaxies and quasars \citep{black_holes_synch, synch_radio_agn}. This emission also dominates the radio continuum from star-forming galaxies, such as our own, at frequencies below $\nu \sim 30$ GHz. The relativistic electrons in nearly all synchrotron sources typically exhibit power-law energy distributions. The energy flux distribution for synchrotron radiation is given by $S_{\nu} \ \propto \ \nu^{\alpha}$, or equivalently, in terms of the brightness temperature as $T_{\nu} \ \propto \ \nu^{\beta}$, where $\beta = \alpha - 2$ \citep{Rybicki_1979}. This frequency dependence of the brightness temperature will be referred to as the spectral index. The spectral index depends on the energy-independent spectral energy slope $s$ of the cosmic-ray electron spectrum, which is typically assumed to have a value of $s=-3$ in the frequency range from 10 to 30 GHz, but varying at lower frequencies \citep{Padovani_2021}. Spatial variations of $\beta$ reflect spatial variations of the cosmic-ray electrons and magnetic field properties in the interstellar medium along the line of sight. Full-sky maps of Galactic synchrotron emission have already been derived in \cite{Krachmalnicoff_2018} or \cite{QUIJOTE_IV}, and their analyses clearly show spatial variations of the spectral index, as discussed below.

\paragraph{}

At higher frequencies, the dominant polarized foreground is thermal emission from interstellar dust grains aligned with the GMF. This emission dominates the microwave sky above approximately 100 GHz and is partially linearly polarized. The frequency spectrum of thermal dust emission is well described by a modified black body (grey body), $I_d(\nu) \propto \nu^{\beta_d} B_\nu(T_d)$, where $B_\nu(T_d)$ is the Planck black-body spectrum at dust temperature $T_d$. Using combined Planck and WMAP data, \citet{Planck_XXII} determined the dust spectral index in both intensity and polarization, finding $\beta_d \simeq 1.59$ for $T_d \simeq 19.6\,\mathrm{K}$, while \citet{Planck_XI} derived a mean polarization spectral index of $\beta_d = 1.53 \pm 0.02$ from PR3 maps at 353 GHz, with the dust $EE$ and $BB$ power spectra following power laws in multipole and exhibiting a significant $E/B$ asymmetry. Since both dust and synchrotron polarization are aligned by the same GMF, a spatial cross-correlation between the two components is expected and has indeed been detected \citep{Choi_2015, Planck_XI}.

\paragraph{}

The sensitivities of WMAP and Planck alone do not allow a sufficiently accurate characterization of the polarized synchrotron signal at intermediate and high Galactic latitudes, leaving important uncertainties in the modelling of this foreground component. This is significantly improved by the QUIJOTE experiment \citep{QUIJOTE_IV}, which observes the polarized Northern sky in the 10--40 GHz range, precisely where synchrotron emission is strongest and where Faraday rotation effects are not expected to be significant \citep{faraday_rotations}. Combining QUIJOTE with WMAP and Planck data, \cite{QUIJOTE_VIII} performed one of the first reconstructions of the polarized synchrotron spectral index in real space, finding statistically significant spatial variations across the sky, with an average value of $\beta = -3.08$ and a dispersion of 0.13. More recently, \cite{casas_2025} studied the spatial variability of polarized synchrotron emission using QUIJOTE MFI data alone, reinforcing the evidence for spatial variability in polarized synchrotron emission and highlighting the importance of combining QUIJOTE observations with complementary data from WMAP, Planck, and future low-frequency surveys in order to reduce noise contamination and improve the robustness of spectral-index estimates.


\paragraph{}

In this work, we characterize the spectral properties of polarized synchrotron and thermal dust emission through power spectrum analyses of the combined QUIJOTE MFI DR1 (11, 13, 17, and 19 GHz), WMAP 9-year (K, Ka, Q, V, and W bands), and Planck PR3 (LFI and HFI, 28.4--353 GHz) datasets. Both foreground components are modelled simultaneously by fitting a parametric SED model to the $EE$ and $BB$ angular power spectra (APS) and their cross-spectra across this wide frequency range. A key aspect of this analysis is the explicit treatment of the spatial cross-correlation between polarized synchrotron and thermal dust emission. Their cross-spectra encode complementary information about the geometry and coherence of the Galactic magnetic field along the line of sight \citep{Planck_XXII, Choi_2015}. The analysis is carried out over the full Northern sky footprint of QUIJOTE, using three Galactic latitude masks to assess the sensitivity of the results to Galactic plane contamination. The resulting foreground characterization is intended to serve as legacy input for component separation in upcoming CMB polarization experiments such as LiteBIRD \citep{LiteBIRD}.